\section{Komplexné čísla a iné polia}

Doteraz sme vektorový priestor chápali ako štvoricu $(V, +, \cdot, \vec{0})$, kde:
\begin{itemize}
    \item $V$ je množina (jej prvky nazývame vektory),
    \item $+: V \times V \to V$ je operácia sčítania vektorov,
    \item $\cdot: \R \times V \to V$ je operácia násobenia vektora skalárom,
    \item $\vec{0} \in V$ je nulový vektor.
\end{itemize}

Samotná existencia tejto štruktúry však na definíciu vektorového priestoru nestačí; vyžaduje sa splnenie určitých vlastností (axióm). Na základe týchto axióm sme následne budovali teóriu vektorových priestorov.

Keďže sme pri budovaní teórie vychádzali výlučne z axióm, všetky zavedené pojmy a dokázané tvrdenia platia pre ľubovoľné vektorové priestory, ktoré tieto axiómy spĺňajú.

Tento koncept teraz zovšeobecníme a zavedieme doň väčšiu flexibilitu.

\begin{tvrdenie}[Pozorovanie]
Všetky tvrdenia dokázané v prvom semestri by platili aj vtedy, ak by sme množinu reálnych čísel $\R$ nahradili inou štruktúrou, ktorej prvky vieme sčítať, odčítať a násobiť. Tieto operácie sme so skalármi nepochybne vykonávali.
\end{tvrdenie}

Aké iné množiny, ktorých prvky vieme sčítať, odčítať a násobiť, poznáme?
\begin{itemize}
    \item Prichádza do úvahy množina prirodzených čísel $\Nat$? Mohli by sme budovať lineárnu algebru nad takýmito skalármi? Nie. Potrebujeme totiž aj deliť (napríklad pri riešení sústav lineárnych rovníc).
    \item Na druhej strane, množina $\R$ má mnoho ďalších vlastností (napríklad usporiadanie, konvergencia postupností), ktoré sme v lineárnej algebre nevyužívali.
\end{itemize}

\begin{question}
Aké vlastnosti sčítania a násobenia v $\R$ boli pre budovanie našej teórie nevyhnutné?
\end{question}
Odpoveďou je nasledujúca definícia.

\begin{definicia}
\textbf{Pole} je pätica $(K, +, \cdot, 0, 1)$, kde:
\begin{itemize}
    \item $K$ je množina,
    \item $+: K \times K \to K$ a $\cdot: K \times K \to K$ sú binárne operácie,
    \item $0 \in K$ a $1 \in K$ sú konštanty,
\end{itemize}
pričom sú splnené nasledujúce axiómy poľa:
\begin{enumerate}
    \item $\forall a,b,c \in K: (a+b)+c = a+(b+c)$ (asociatívnosť sčítania)
    \item $\forall a,b \in K: a+b = b+a$ (komutatívnosť sčítania)
    \item $\forall a \in K: a+0 = a$ (existencia neutrálneho prvku pre sčítanie)
    \item $\forall a \in K, \exists b \in K: a+b=0$ (toto $b$ je nutne jediné a značí sa $-a$, existencia opačného prvku)
    \item $\forall a,b,c \in K: (a \cdot b) \cdot c = a \cdot (b \cdot c)$ (asociatívnosť násobenia)
    \item $\forall a \in K: 1 \cdot a = a$ (existencia neutrálneho prvku pre násobenie)
    \item $\forall a \in K, a \neq 0, \exists b \in K: a \cdot b = 1$ (toto $b$ je nutne jediné a značí sa $a^{-1}$, existencia inverzného prvku)
    \item $\forall a,b,c \in K: a \cdot (b+c) = (a \cdot b) + (a \cdot c)$ (distributívnosť)
    \item $0 \neq 1$ (netriviálnosť)
\end{enumerate}
\end{definicia}

\subsection{Príklady polí}
Dva príklady sú nám dobre známe: reálne čísla $\R$ a racionálne čísla $\mathbb{Q}$. Väčšinu času tejto prednášky venujeme zavedeniu a skúmaniu ďalšieho príkladu: komplexných čísel $\mathbb{C}$.

Skôr než začneme, spomeňme ešte jednu triedu príkladov: konečné polia $Z_p$.

Nech $p$ je prvočíslo. Na množine $Z_p = \{0, 1, \dots, p-1\}$ definujme operácie sčítania a násobenia nasledovne:
\begin{itemize}
    \item $a \oplus b = (a+b) \pmod p$
    \item $a \odot b = (a \cdot b) \pmod p$
\end{itemize}

Napríklad pre $p=3$ dostávame tabuľky operácií:

\begin{center}
\begin{tabular}{c|ccc}
$\oplus$ & 0 & 1 & 2 \\
\hline
0 & 0 & 1 & 2 \\
1 & 1 & 2 & 0 \\
2 & 2 & 0 & 1
\end{tabular}
\hspace{2cm}
\begin{tabular}{c|ccc}
$\odot$ & 0 & 1 & 2 \\
\hline
0 & 0 & 0 & 0 \\
1 & 0 & 1 & 2 \\
2 & 0 & 2 & 1
\end{tabular}
\end{center}

Potom štruktúra $(Z_p, \oplus, \odot, 0, 1)$ tvorí pole (dôkaz vynechávame). Poznamenajme, že ak $p$ nie je prvočíslo, axiómy poľa nie sú splnené.
Polia $Z_p$ sú veľmi užitočné v aplikáciách (teória kódovania, kryptografia, \dots), my sa nimi však hlbšie zaoberať nebudeme.

\subsection{Komplexné čísla}

Cardanov vzorec z roku 1545 na riešenie kubickej rovnice obsahuje odmocniny zo záporných čísel, a to aj v prípade, keď sú všetky riešenia reálne. Použitie tohto vzorca teda predpokladalo existenciu odmocniny zo záporného čísla.

Uvažujme reálne čísla $\R$. Pridajme k nim nový prvok $i \notin \R$ a definujme $i^2 = -1$. Množinu postupne rozširujme tak, aby sme pridali len nevyhnutné prvky na vytvorenie poľa. Pribudnú prvky tvaru $bi$ a $a+bi$, kde $a, b \in \R$. Tento proces je postačujúci.

\begin{definicia}[Komplexné čísla]
Množinu komplexných čísel definujeme ako $\mathbb{C} = \R \times \R$. Prvky nezapisujeme ako usporiadané dvojice $(a, b)$, ale v tvare:
\[ a + bi \]
kde $a, b \in \R$. Tento zápis sa nazýva \textbf{algebraický tvar}.
Pritom stotožňujeme 
\begin{itemize}
    \item $\R\ni a \equiv (a, 0)$
    \item $bi \equiv (0, b)$
\end{itemize}
Sčítanie $+: \mathbb{C} \times \mathbb{C} \to \mathbb{C}$ je definované predpisom:
\[ (a_1 + b_1 i) + (a_2 + b_2 i) = (a_1 + a_2) + (b_1 + b_2)i \]
Násobenie $\cdot: \mathbb{C} \times \mathbb{C} \to \mathbb{C}$ je definované predpisom:
\[ (a_1 + b_1 i)(a_2 + b_2 i) = (a_1 a_2 - b_1 b_2) + (a_1 b_2 + b_1 a_2)i \]
(Toto zodpovedá pravidlám pre násobenie dvojíc).
Nulový prvok je $0 = 0 + 0i$. Jednotkový prvok je $1 = 1 + 0i$.
Štruktúra $( \mathbb{C}, +, \cdot, 0, 1 )$ potom tvorí pole.
\end{definicia}

V zmysle stotožnenia v predošlej definícii stotožňujeme reálne číslo $a$ s komplexným číslom $a+0i$.

\begin{veta}
Komplexné čísla tvoria pole.
\end{veta}
\begin{proof}
(iba náčrt) Axiómy pre sčítanie sú zrejmé. Pre násobenie: ak $z = a+bi$ a $z \neq 0$, ako vyzerá $z^{-1}$?
\[ z^{-1} = \frac{1}{z} = \frac{1}{a+bi} \]
Zlomok rozšírime výrazom $a-bi$:
\[ \frac{1}{a+bi} \cdot \frac{a-bi}{a-bi} = \frac{a-bi}{a^2 - abi + abi + b^2} = \frac{a-bi}{a^2+b^2} = \frac{a}{a^2+b^2} - \frac{b}{a^2+b^2}i \]
Teda
\[ (a+bi)^{-1} = \frac{a}{a^2+b^2} - \frac{b}{a^2+b^2}i \]
Je potrebné overiť, že platí: $(a+bi) \cdot z^{-1} = 1$, to môžete urobiť sami.

Ostatné vlastnosti operácií sčítania a násobenia komplexných čísel je možné overiť priamym výpočtom (čo môže byť prácne,
ale jednoduché, preto to nebudeme robiť).
\end{proof}

\subsection{Geometrická interpretácia}
Keďže komplexné číslo $a+bi$ možno stotožniť s usporiadanou dvojicou čísel $(a,b)$, je prirodzené reprezentovať ho ako bod v rovine vybavenej dvoma kolmými súradnicovými osami:
\begin{center}
\begin{tikzpicture}[>=Stealth, thick]

    % Define styles
    \tikzset{
        graphline/.style={blue!80!black},
        guide/.style={blue!80!black, densely dotted},
        % Standard brace (appears on 'left' side of the drawing path)
        brace/.style={decorate, decoration={brace, amplitude=6pt, raise=2pt}, blue!80!black},
        % Mirrored brace (appears on 'right' side of the drawing path)
        brace mirror/.style={decorate, decoration={brace, mirror, amplitude=6pt, raise=2pt}, blue!80!black},
        labeltext/.style={blue!80!black},
        rednote/.style={red!85!black, font=\small\itshape},
        redarrow/.style={->, red!85!black, shorten >=2pt, shorten <=2pt}
    }

    % --- LEFT DIAGRAM ---
    \begin{scope}
        % Axes
        \draw[graphline] (-1.5,0) -- (3.5,0);
        \draw[graphline] (0,-2.5) -- (0,3);

        % Coordinates
        \def\xa{2}
        \def\yb{1.5}
        \coordinate (P1) at (\xa, \yb);

        % Dotted lines
        \draw[guide] (P1) -- (\xa, 0);
        \draw[guide] (P1) -- (0, \yb);

        % Point
        \fill[graphline] (P1) circle (1.5pt) node[anchor=west, xshift=3pt] {$a+bi$};

        % Braces and Labels
        
        % Horizontal brace (a) - Below axis
        % Path: Left -> Right. Mirror puts it on the Right side (Below).
        \draw[brace mirror] (0,0) -- (\xa,0) node[midway, below=10pt, labeltext] {$a$};
        
        % Vertical brace (b) - Left of axis
        % Path: Bottom -> Top. Standard brace puts it on the Left side.
        \draw[brace] (0,0) -- (0,\yb) node[midway, left=10pt, labeltext] {$b$};
    \end{scope}


    % --- RIGHT DIAGRAM ---
    \begin{scope}[xshift=8cm]
        % Axes
        \draw[graphline] (-2.5,0) -- (3.5,0);
        \draw[graphline] (0,-2.5) -- (0,3);

        % Coordinates
        \def\xa{2}
        \def\yb{1.5}
        \coordinate (P2) at (\xa, \yb);

        % Dotted lines
        \draw[guide] (P2) -- (\xa, 0);
        \draw[guide] (P2) -- (0, \yb);

        % Point
        \fill[graphline] (P2) circle (1.5pt) node[anchor=west, xshift=3pt] {$z$};

        % Braces and Labels
        
        % Horizontal brace (Re z) - Below axis
        \draw[brace mirror] (0,0) -- (\xa,0) node[midway, below=10pt, labeltext] (ReZ) {$\Re z$};
        
        % Vertical brace (Im z) - Left of axis
        % Path: Bottom -> Top. Standard brace puts it on the Left side.
        \draw[brace] (0,0) -- (0,\yb) node[midway, left=10pt, labeltext] (ImZ) {$\Im z$};

        % --- Red Annotations ---
        
        % Imaginarna cast label (Top Left)
        % Adjusted position to account for brace being on the left now
        \node[rednote, anchor=east] (ImText) at (-1.0, 2.5) {imaginárna časť $z$};
        % Arrow pointing from Im z label to the text
        \draw[redarrow] (ImZ.west) to (ImText.south);

        % Realna cast label (Bottom)
        \node[rednote, anchor=north] (ReText) at (\xa, -1.8) {reálna časť $z$};
        % Arrow pointing from Re z label to the text
        \draw[redarrow] (ReZ.south) -- (ReText.north);

    \end{scope}

\end{tikzpicture}
\end{center}
\begin{itemize}
    \item Reálna časť $z$ ($\Re z$) sa vynáša na os x.
    \item Imaginárna časť $z$ ($\Im z$) sa vynáša na os y.
\end{itemize}

Sčítanie komplexných čísel má potom zrejmý geometrický význam, analogický sčítaniu vektorov v rovine.
Aj násobenie komplexných čísel možno interpretovať geometricky. Skôr ako to urobíme,
zavedieme si niekoľko pojmov.

\begin{definicia}~
\begin{itemize}
    \item \textbf{Absolútna hodnota} komplexného čísla $z = a+ib$ je $|z| = \sqrt{a^2+b^2}$.
    \item \textbf{Konjugované číslo} (združené) k číslu $z= a+ib$ je $\overline{z} = a-ib = \Re z - i(\Im z)$.
\end{itemize}
\end{definicia}


Platí $z \cdot \overline{z} = (a+bi)(a-bi) = a^2+b^2 \in \R$. Teda $|z|^2 = z \cdot \overline{z}$, resp. $|z| = \sqrt{z \cdot \overline{z}}$.

Všimnime si, že vyššie odvodený vzorec pre inverzný prvok možno zapísať ako $z^{-1} = \frac{\overline{z}}{|z|^2}$.

Geometricky je \[
|z|=\sqrt{(\Re z)^2+(\Im z)^2}
\] vzdialenosť od počiatku súradnicovej sústavy. 

\subsection{Polárny tvar}

Vieme, že ak uvažujeme rovinu s počiatkomi $O$, zavedenie dvoch kolmých súradnicových osí pretínajúcich sa v počiatku s
rovnakou vzdialenosťou bodu $1$ od bodu $0$ nám umožňuje reprezentovať body v rovine ako usporiadané dvojice reálnych
čísel. Ale v rovine môžeme reprezentovať body aj inak, bod $A$ vieme jednoznačne reprezentovať pomocou dvojice
\[
(\text{dĺžka $AO$},\text{uhol zovretý polpriamkou $OA$ a kladnou polosou $x$})
\]
v prípade, že $A=O$, uhol zvolíme rovný $0$, toto nespôsobuje ťažkosti. Táto dvojica sa nazýva \emph{polárne súradnice}
bodu $A$ v rovine.

Ak $z\in\mathbb C$, vzdialenosť $z$ od počiatku komplexnej roviny je $|z|$. Uhol nazveme \emph{argument $z$}. Celá
geometrická situácia vyzerá teraz takto
\begin{center}
\begin{tikzpicture}[>=Stealth, thick]

    % Define styles
    \tikzset{
        axisline/.style={blue!80!black},
        guide/.style={blue!80!black, densely dotted},
        redpath/.style={red!85!black, thick},
        redarrow/.style={->, red!85!black},
        greenarrow/.style={->, green!60!black, shorten >=2pt, shorten <=2pt},
        bluetext/.style={blue!80!black},
        redtext/.style={red!85!black},
        greentext/.style={green!60!black, font=\small\itshape, align=center},
        brace/.style={decorate, decoration={brace, amplitude=6pt, raise=4pt}, red!85!black},
        brace mirror/.style={decorate, decoration={brace, mirror, amplitude=6pt, raise=4pt}, red!85!black}
    }

    % 1. Define Coordinates
    \coordinate (O) at (0,0);
    \def\xa{3}
    \def\yb{2}
    
    \coordinate (Z) at (\xa, \yb);          % Point z
    \coordinate (Zbar) at (\xa, -\yb);      % Point z_bar (conjugate)
    \coordinate (XPoint) at (\xa, 0);       % Projection on x-axis

    % 2. Axes
    \draw[axisline] (-2.5,0) -- (5.5,0);
    \draw[axisline] (0,-3.5) -- (0,3.5);

    % 3. Z Elements (Upper Half)
    \draw[guide] (Z) -- (XPoint);
    \draw[guide] (Z) -- (0, \yb);
    \draw[redpath] (O) -- (Z);
    \node[bluetext, right, xshift=2pt] at (Z) {$z$};

    % 4. Z_bar Elements (Lower Half - New)
    % Guide lines for z_bar
    \draw[guide] (Zbar) -- (XPoint);
    \draw[guide] (Zbar) -- (0, -\yb);
    % Vector mirroring z
    %\draw[redpath] (O) -- (Zbar);
    % Node label
    \node[bluetext, right, xshift=2pt] at (Zbar) {$\bar{z}$};
    
    % Formula for z_bar placed near the node

    % 5. Braces (Red)
    % Re z (horizontal, mirror to be below)
    \draw[brace mirror] (O) -- (XPoint) node[midway, below=12pt, redtext] (ReZ) {$\Re z$};

    % Im z (vertical, standard to be on left)
    \draw[brace] (0,0) -- (0,\yb) node[midway, left=12pt, redtext] (ImZ) {$\Im z$};

    % |z| (diagonal along vector)
    \draw[brace, decoration={raise=2pt}] (O) -- (Z) node[midway, above left=8pt, redtext, sloped] (AbsZ) {$|z|$};

    % 6. Angle (Argument)
    \pic [draw=red!85!black, ->, angle radius=1.0cm] {angle = XPoint--O--Z};
    
    % Label for argument with arrow
    \node[redtext] (ArgLabel) at (2.5, 0.8) {$\arg(z)$};
    \draw[redarrow, bend right=10] (ArgLabel.west) to (1.0, 0.3);

    % 7. Green Annotations
    % Absolutna hodnota
    \node[greentext, anchor=west] (AbsText) at (1.5, 3.2) {absolútna hodnota $z$};
    \draw[greenarrow] (AbsZ.north) to [bend left=10] (AbsText.south west);

    % Argument z
    \node[greentext, anchor=west] (ArgText) at (4.5, 1.8) {argument $z$};
    \draw[greenarrow] (ArgLabel.east) -- (ArgText.west);

    % Imaginarna cast
    \node[greentext, anchor=south east] (ImText) at (-1, 2.5) {imaginárna\\časť $z$};
    \draw[greenarrow] (ImZ.west) to [bend left=10] (ImText.south);

    % Realna cast
    \node[greentext, anchor=north] (ReText) at (4.0, -2.8) {reálna časť $z$};
    \draw[greenarrow] (ReZ.south) to [bend right=10] (ReText.north);
    
    % Absolutna hodnota
    \node[greentext, anchor=west,xshift=40pt,yshift=40pt] (ZbarText) at (Zbar) {komplexne združené k $z$};
    \draw[greenarrow] (Zbar.east)+(0.4,0.2) to (ZbarText.west);

\end{tikzpicture}
\end{center}

Z pravouhlého trojuholníka v komplexnej rovine dostaneme:
\[ \Im z = |z| \sin(\arg z) \]
\[ \Re z = |z| \cos(\arg z) \]
kde $\arg z$ je uhol, ktorý zviera sprievodič bodu $z$ s kladnou polosou x.
Dosadením dostávame 
\[
z = \Re z + i(\Im z)= |z|\cos(\arg z) + i|z|\sin(\arg z)
=|z|(\cos(\arg z)+i\sin(\arg z))
\]
Teda

Analyzujme násobenie komplexných čísel v goniometrickom tvare. Nech $z = |z|(\cos \alpha + i \sin \alpha)$ a $w = |w|(\cos \beta + i \sin \beta)$.
\[ z \cdot w = |z||w| (\cos \alpha + i \sin \alpha)(\cos \beta + i \sin \beta) \]
\[ = |z||w| [ (\cos \alpha \cos \beta - \sin \alpha \sin \beta) + i (\sin \alpha \cos \beta + \cos \alpha \sin \beta) ] \]
Použitím súčtových vzorcov:
\[ \cos(\alpha+\beta) = \cos \alpha \cos \beta - \sin \alpha \sin \beta \]
\[ \sin(\alpha+\beta) = \sin \alpha \cos \beta + \cos \alpha \sin \beta \]
dostávame:
\[ z \cdot w = |z||w| (\cos(\alpha+\beta) + i \sin(\alpha+\beta)) \]

Tento vzorec má nasledujúcu geometrickú interpretáciu: pri násobení komplexných čísel sa \textbf{sčítajú ich argumenty} a \textbf{násobia ich absolútne hodnoty}.
\[ |zw| = |z||w| \]
\[ \arg(zw) = \arg z + \arg w \]

\begin{center}
\begin{tikzpicture}[>=Stealth, thick, line cap=round]

    % --- Definitions ---
    % Colors to match the image
    \definecolor{myred}{RGB}{220, 20, 60}
    \definecolor{myblue}{RGB}{25, 25, 180}
    \definecolor{mygreen}{RGB}{0, 150, 60}

    % Coordinates (Polar: angle:radius)
    \def\angZ{35}   % Red angle
    \def\lenZ{3.5}
    
    \def\angW{55}   % Blue angle
    \def\lenW{4.5}
    
    \def\angZW{110} % Green angle = sum
    \def\lenZW{5.5} % Schematic length for z*w

    \coordinate (O) at (0,0);
    \coordinate (X) at (4,0); % Reference for angles
    
    \coordinate (Z) at (\angZ:\lenZ);
    \coordinate (W) at (\angW:\lenW);
    \coordinate (ZW) at (\angZW:\lenZW);

    % --- Axes ---
    \draw[black] (-2.5,0) -- (4.0,0); % x-axis
    \draw[black] (0,-1.0) -- (0,6); % y-axis

    % --- Vectors and Points ---
    
    % Red Vector (z)
    \draw[, myred] (O) -- (Z);
    \fill[myred] (Z) circle (2pt);
    \node[myred, anchor=south west] at (Z) {$z$};

    % Blue Vector (w)
    \draw[, myblue] (O) -- (W);
    \fill[myblue] (W) circle (2pt);
    \node[myblue, anchor=south west, xshift=2pt, yshift=2pt] at (W) {$w$};

    % Green Vector (z.w)
    \draw[, mygreen] (O) -- (ZW);
    \fill[mygreen] (ZW) circle (2pt);
    \node[mygreen, anchor=south, yshift=2pt] at (ZW) {$z \cdot w$};
    
    % --- Green Bracket and Label ---
    % Label is now horizontal (no 'sloped') and positioned to the left
    \draw[decorate, decoration={brace, amplitude=6pt, raise=5pt}, mygreen] 
        (O) -- (ZW) 
        node[midway, left=15pt, mygreen] {$|z||w|$};

    % --- Angles and Labels ---

    % 1. Red Angle (arg z)
    \draw[myred] (1.8,0) arc (0:\angZ:1.8);
    \node[myred, anchor=west] at (15:2.0) {$\arg z$};

    % 2. Blue Angle (arg w)
    \draw[myblue] (1.2,0) arc (0:\angW:1.2);
    \node[myblue, rotate=\angW/2+15, anchor=west] at (\angW/2+15:1.4) {$\arg w$};

    % 3. Green Angle (arg(z) + arg(w))
    \draw[mygreen] (0.7,0) arc (0:\angZW:0.7);
    \node[mygreen, rotate=\angZW/2+10, anchor=south west] at (\angZW/2:0.8) {$\arg(z)+\arg(w)$};

\end{tikzpicture}
\end{center}

\begin{priklad}
Nech $z = 1+i$, $w = \sqrt{3}+i$.
Platí $|z| = \sqrt{2}$, $\arg z = \frac{\pi}{4}$.
Ďalej $|w| = 2$, $\arg w = \frac{\pi}{6}$.
Súčin $z \cdot w$ má potom absolútnu hodnotu $2\sqrt{2}$ a argument $\frac{\pi}{4} + \frac{\pi}{6} = \frac{5\pi}{12}$.
\end{priklad}


\subsubsection{Moivreova veta}
Induktívnym rozšírením vzorca pre násobenie dostaneme vzorec pre mocninu komplexného čísla. Pre $n \in \Nat$ platí:
\[ z^n = |z|^n (\cos(n \arg z) + i \sin(n \arg z)) \]

\subsubsection{Riešenie rovnice $z^n = c$}
Predpokladajme, že máme $n \in \Nat$ a číslo $c \in \mathbb{C}$. Našou úlohou je vyriešiť rovnicu $z^n = c$ (t.j. nájsť $n$-tú odmocninu z $c$).
Vyjadrime obe strany rovnice v goniometrickom tvare. Nech $z = |z|(\cos \varphi + i \sin \varphi)$ a $c = |c|(\cos \psi + i \sin \psi)$.
Rovnica $z^n = c$ má potom tvar:
\[ |z|^n (\cos n\varphi + i \sin n\varphi) = |c|(\cos \psi + i \sin \psi) \]
Porovnaním absolútnych hodnôt dostávame $|z|^n = |c|$, teda $|z| = \sqrt[n]{|c|}$.
Pre argumenty platí:
\[ \cos n\varphi = \cos \psi \]
\[ \sin n\varphi = \sin \psi \]
Aké sú možné hodnoty $n\varphi$?
Jedným z riešení je $\psi$, ale musíme zohľadniť periodicitu funkcií sínus a kosínus (perióda $2\pi$). Všeobecné riešenie má tvar:
\[ n\varphi = \psi + 2k\pi, \quad k \in \{0, 1, \dots, n-1\} \]
Teda
\[ \varphi_k = \frac{\psi + 2k\pi}{n} \]

Geometricky tvoria riešenia vrcholy pravidelného $n$-uholníka ležiaceho na kružnici s polomerom $\sqrt[n]{|c|}$.

\begin{priklad}
Riešme rovnicu $z^3 = -2 + 2i$.
Máme $c = -2+2i$. Absolútna hodnota je $|c| = \sqrt{4+4} = \sqrt{8}$.
Argument je $\arg c = \frac{3\pi}{4}$ (pretože reálna časť je záporná, imaginárna kladná a sú v absolútnej hodnote rovnaké).
Riešenia ležia na kružnici so stredom v 0 a polomerom $\sqrt[3]{\sqrt{8}} = \sqrt{2}$.
Argumenty riešení sú:
\[ \varphi_k = \frac{\frac{3\pi}{4} + 2k\pi}{3} = \frac{\pi}{4} + \frac{2k\pi}{3} \]
Pre $k=0$ máme $\varphi_0 = \frac{\pi}{4}$. Prvé riešenie je $z_0 = \sqrt{2}(\cos \frac{\pi}{4} + i \sin \frac{\pi}{4}) = \sqrt{2}(\frac{\sqrt{2}}{2} + i\frac{\sqrt{2}}{2}) = 1+i$.
Pre $k=1$ a $k=2$ dostaneme ostatné vrcholy trojuholníka.
\end{priklad}

